\subsection{The finite-flat section sandwich and the actual norm degree}
\label{subsec:faithful-geometric-norm-degree}

This subsection proves the geometric part of
\srcid{DegreeMultiplicationNorm}.  In particular it includes
inseparable multiplication maps: the section comparison is for
finite flat covers, with no \'etaleness assumption.

\paragraph{Exact classical background for the section input.}
The coherent background used is Serre's theorem on projective
space over a Noetherian ring: coherent cohomology is finite,
and positive cohomology of sufficiently high twists vanishes
(\href{https://stacks.math.columbia.edu/tag/01YS}{Stacks, Lemma 30.14.1}).
In our application the base is a field \(k\), and the specified
very ample sections define a closed immersion
\(\iota:X\hookrightarrow\mathbb P^m_k\) with
\(H=\iota^*\mathcal O(1)\).
Indeed the source chart data say that the nonzero loci cover
\(X\), are affine, and have coordinate algebras generated by
the indicated section ratios.  The induced maps on all
standard projective charts are closed immersions.
The ideal sheaf \(\mathcal I_X\) is coherent.  Twisting
\(0\to\mathcal I_X\to\mathcal O\to\iota_*\mathcal O_X\to0\)
and applying the vanishing of \(H^1(\mathcal I_X(t))\)
shows that, for all sufficiently large \(t\), every section
of \(H^t\) is a homogeneous polynomial of degree \(t\) in those
specified sections.  This is precisely the needed
\srcid{ExistsHomogeneousFormEqHighPowerSection} interface.
There is no extra uniformity in \(X\), its embedding, or field
extension in that assertion.

We also use the ample-sheaf theorem that a finite-type
quasi-coherent sheaf tensored with sufficiently high powers
of an ample line is globally generated
(\href{https://stacks.math.columbia.edu/tag/01Q3}{Stacks, Proposition 28.27.13(7)}).
For a line \(N\) on a nonempty projective \(X\), it gives a
nonzero section of \(N\otimes H^a\) for some \(a\).
This matches the source's
\srcid{ExistsHighTwistSectionNeZero}, on its actual field
and projective scheme.
Both uses have Noetherian base \(k\); neither asserts a result
over an arbitrary non-Noetherian local model.

\begin{srlemma}[The source section growth]\label{lem:faithful-section-growth}
Let \(X/k\) be integral and smooth of relative dimension \(g\),
with specified very ample line \(H\) and sections
\(s_0,\ldots,s_m\).
Then \(h(t)=\dim_kH^0(X,H^t)\) is finite for every \(t\ge0\),
is nondecreasing, and there are integers \(c,C\), \(c\ge1\), with
\[
                    t^g\le c h(t),\qquad h(t)\le Ct^g
                        \quad(t\ge1).
 \tag{DG1}\label{eq:faithful-section-growth}
\]
\end{srlemma}
\begin{proof}
We give the source's filtered-coordinate proof.
For a finitely generated integral \(k\)-algebra \(A\) of
transcendence degree \(g\), generated by \(r_1,\ldots,r_m\),
let \(F_tA\) be the image of polynomials of total degree
at most \(t\) in these generators.
Noether normalization supplies algebraically independent
\(y_1,\ldots,y_g\) such that \(A\) is finite over
\(k[y_1,\ldots,y_g]\); choose module generators \(b_1,\ldots,b_l\).
Let \(B_D\) be the span of \(y^\alpha b_j\) with
\(0\le\alpha_i\le D\).
These spaces exhaust \(A\) and have dimension at most
\(l(D+1)^g\).
Choose one \(d\) with all \(r_i b_j\in B_d\), and choose
\(D_0\) with \(1\in B_{D_0}\).
Multiplying \(B_D\) by any \(r_i\) lands in \(B_{D+d}\).
Hence \(F_tA\subset B_{td+D_0}\), giving
\(\dim F_tA\le C_A t^g\) for \(t\ge1\).

For the reverse inequality, lift each \(y_i\) to a polynomial
\(q_i(r)\), and put \(E=1+\sum_i\deg q_i\).
The independent monomials \(y^\alpha\), with
\(0\le\alpha_i\le\lfloor t/E\rfloor\), belong to \(F_tA\).
Thus
\[
             t^g\le E^g\dim_kF_tA .
\]
These are the actual box-span upper and lower arguments in
\srcid{FinrankTotalDegreeImageGrowth}; no Hilbert-polynomial
leading coefficient is substituted.

Now choose a nonempty standard open \(D_i=(s_i\ne0)\).
It is affine, with coordinate algebra generated by
\(r_j=s_j/s_i\), and \(r_i=1\); this is part of the specified
very-ample chart data.
Its transcendence degree is \(g\), since it is a nonempty
affine in the given smooth integral \(g\)-dimensional scheme.
The map
\[
      H^0(X,H^t)\longrightarrow\Gamma(D_i,\mathcal O),
                         \qquad \sigma\longmapsto\sigma/s_i^t
\]
is injective by integrality.
Every polynomial in the \(r_j\) of degree at most \(t\)
is the image of a global section, by homogenizing with the
variable \(r_i=1\).
For all sufficiently large \(t\), the exact homogeneous-section
input stated above gives the reverse inclusion.
Thus, in large degrees, this section space has exactly the
filtered-coordinate dimension just bounded.

Multiplication by the fixed nonzero section \(s_i\) is injective
from \(H^0(H^t)\) to \(H^0(H^{t+1})\), by integrality.
It gives monotonicity and propagates finiteness from the large
degrees to every smaller degree.
Increase \(C\) by the finite sum of these smaller dimensions.
The filtered-coordinate lower bound holds in every degree
because the homogenized polynomials always give global sections.
This proves \eqref{eq:faithful-section-growth}.
\end{proof}

\begin{srlemma}[The actual finite-flat section sandwich]
\label{lem:faithful-finite-flat-section-sandwich}
Let \(X,X'\) be integral \(k\)-schemes, with \(X\) projective,
and let \(\varphi:X'\to X\) be finite, flat, and surjective of
constant rank \(\delta\).
Let \(H\) have the specified finite very-ample chart data on \(X\).
Then some \(a\ge0\) satisfies, for every \(t\ge0\),
\[
 \begin{split}
  \delta\,h^0(X,H^t)
       &\le h^0(X',\varphi^*H^{t+a}),\\
  h^0(X',\varphi^*H^t)
       &\le \delta\,h^0(X,H^{t+a}).
 \end{split}
 \tag{DG2}\label{eq:faithful-finite-flat-sandwich}
\]
\end{srlemma}
\begin{proof}
All the section spaces are finite by coherent finiteness on
the projective schemes \(X\) and \(X'\).
Put \(F=k(X)\), \(F'=k(X')\).  Finite flatness identifies
\([F':F]=\delta\); choose an \(F\)-basis \(e_1,\ldots,e_\delta\)
of \(F'\).
For every nonempty standard open \(D_i=(s_i\ne0)\), put
\[
                 A_i=\Gamma(D_i,\mathcal O_X),\qquad
                 B_i=\Gamma(\varphi^{-1}D_i,\mathcal O_{X'}).
\]
They are domains embedded in \(F,F'\), and \(B_i\) is finite
over \(A_i\).
Choose \(0\ne c_i\in A_i\) clearing both sets of denominators:
\[
 \begin{split}
          c_i\,\operatorname{coord}_{e_j}(B_i)&\subset A_i,\\
                              c_i e_j&\in B_i .
 \end{split}
 \tag{DG3}\label{eq:faithful-chart-denominators}
\]
For the first assertion, clear the finitely many coordinates
of finite \(A_i\)-module generators of \(B_i\).
For the second, write each \(e_j\) as a fraction in \(B_i\);
its denominator is integral over \(A_i\), so it divides a
nonzero element of \(A_i\).
Multiply these finitely many denominators.
This is the source's \srcid{exists\_denominator}.

Fix one rational trivialization of \(H\), and write its
sections \(s_i\) as \(0\ne\rho_i\in F\).
Then \(\rho_i/\rho_k\in A_k\).
On \(D_i\cap D_k\), the function \(c_i\) is regular;
localization at the function \(\rho_i/\rho_k\) on \(D_k\)
therefore gives an integer \(m_{ik}\) with
\[
                  c_i(\rho_i/\rho_k)^{m_{ik}}\in A_k.
\]
Let \(M_i=\sum_k m_{ik}\), let \(a=\sum_iM_i\), and set
\[
                           u=\prod_i c_i\rho_i^{M_i}\ne0.
\]
For each \(k\),
\[
 u=\rho_k^a c_k b_k,\qquad
       b_k=\prod_{i\ne k}c_i(\rho_i/\rho_k)^{M_i}\in A_k.
 \tag{DG4}\label{eq:faithful-common-multiple}
\]
Increasing an exponent from \(m_{ik}\) to \(M_i\) preserves
membership because \(\rho_i/\rho_k\in A_k\).
This is the explicit \srcid{exists\_common\_multiple}
construction.

Embed sections in their function fields using the chosen
rational trivialization.  The exact gluing criterion is
\[
 \begin{split}
 V_t&=\{z\in F:z/\rho_i^t\in A_i\text{ for all }i\}
                         \simeq H^0(X,H^t),\\
 W_t&=\{z'\in F':z'/\rho_i^t\in B_i\text{ for all }i\}
                         \simeq H^0(X',\varphi^*H^t).
 \end{split}
\]
These are equalities of images of actual global sections:
regular local coefficients glue by their common function-field
value, and restriction is injective on an integral scheme.
Now the two linear maps are
\[
 \begin{aligned}
 V_t^\delta&\longrightarrow W_{t+a},
                       &(z_j)&\longmapsto u\sum_j z_j e_j,\\
 W_t&\longrightarrow V_{t+a}^\delta,
                       &z'&\longmapsto
                              (u\,\operatorname{coord}_{e_j}(z'))_j .
 \end{aligned}
 \tag{DG5}\label{eq:faithful-sandwich-maps}
\]
They land in the indicated section images by
\eqref{eq:faithful-chart-denominators} and
\eqref{eq:faithful-common-multiple}.
For example the first expression divided by \(\rho_i^{t+a}\)
is \(\sum_j b_i(z_j/\rho_i^t)c_i e_j\in B_i\);
the second uses \(b_i c_i\operatorname{coord}_{e_j}
(z'/\rho_i^t)\in A_i\).
They are injective because \(u\ne0\) and the \(e_j\) are a
basis.  Their dimensions prove
\eqref{eq:faithful-finite-flat-sandwich}.
\end{proof}

\begin{srlemma}[The source growth comparison]
\label{lem:faithful-degree-growth-comparison}
Suppose a nondecreasing function \(h:\mathbb N\to\mathbb N\)
satisfies \eqref{eq:faithful-section-growth}, and \(q\ge1\).
If
\[
         \delta h(t)\le h(q(t+a)),\qquad
                 h(qt)\le\delta h(t+a)\quad(t\ge0),
\]
then \(\delta=q^g\).
\end{srlemma}
\begin{proof}
The lower bound at \(1\) gives \(h(1)\ge1\).
Induction, using monotonicity and \(q\ge1\), gives
\[
 \begin{split}
 h(q^kt)&\le\delta^k h(t+ka),\\
 \delta^k&\le h(q^k(1+ka)).
 \end{split}
\]
Applying the two bounds in
\eqref{eq:faithful-section-growth} yields
\[
 \delta^k\le C(q^g)^k(1+ka)^g,\qquad
       (q^g)^k\le cC\,\delta^k(1+ka)^g .
\]
An exponential with ratio greater than \(1\) eventually
exceeds any fixed polynomial.
The first inequality gives \(\delta\le q^g\);
the second gives \(q^g\le\delta\).
This is the source's \srcid{degree\_eq\_pow\_of\_growth}
and its \srcid{le\_of\_pow\_le\_poly\_mul\_pow} argument.
\end{proof}

\begin{srproposition}[The actual nonzero multiplication kernel]
\label{prop:faithful-proper-multiplication-degree}
Let \(k\) be algebraically closed, \(B/k\) an abelian variety
of dimension \(g\), \(E\) a number field, and
\(\rho:O=\mathcal O_E\to\operatorname{End}(B)\) a unital ring
action.
For \(0\ne a\in O\), the map \([a]\) is finite, flat, and
surjective, and its finite kernel has coordinate algebra
\(A_a\) with
\[
                    (\dim_kA_a)^{[E:\mathbb Q]}
                             =|N_{E/\mathbb Q}(a)|^{2g}.
 \tag{DG6}\label{eq:faithful-proper-multiplication-norm-degree}
\]
This is valid in positive characteristic when the map is
inseparable.
\end{srproposition}
\begin{proof}
An abelian variety is projective; the theorem of the cube
holds for each of its lines
(\href{https://stacks.math.columbia.edu/tag/0BF9}{Stacks,
Lemmas 39.9.2 and 39.9.6}).
Choose a symmetric very ample line \(H\): symmetrize an
ample line \(H_0\) to \(H_0\otimes[-1]^*H_0\), and
take a sufficiently high power.
Write \(h=[H]\) in its Picard group and set
\[
                         Q(a)=[a]^*h .
\]
The cube identity and symmetry give
\[
 Q(0)=0,\qquad Q(a+b)+Q(a-b)=2Q(a)+2Q(b),\qquad
                         [m]^*H\simeq H^{m^2}.
 \tag{DG7}\label{eq:faithful-abelian-quadratic-pull}
\]
The last identity follows by the two-step quadratic induction,
including integers divisible by the characteristic.

We first explain the source's finiteness criterion.
If a morphism from a projective scheme with very ample line
\(H'\) pulls a line on its separated target back to
\((H')^D\), \(D\ge1\), it is finite.
It is proper.  On each fibre the pulled-back line is a
constant line, so its trivialization makes the \(D\)-th
powers of the very ample sections into global functions.
Their basic opens are exactly the inverse images of the
standard very-ample affine opens, and are affine.
The fibre is proper over its residue field, so its global
functions form a finite-dimensional algebra.
Localizing that finite algebra at each of those functions
shows that these affine basic opens are finite over the
residue field.  They cover the fibre, so it is zero-dimensional.
The morphism is locally quasi-finite and proper, hence finite.
This is \srcid{IsFiniteOfPullVeryAmpleIsoPow}, including its
construction of the fibre functions from local line frames.

Apply this with \eqref{eq:faithful-abelian-quadratic-pull}
to every positive integer multiplication.
For nonzero \(a\), the integer \(M=|N(a)|\) lies in \((a)\);
write \(M=ab\).
Then \([M]=[b][a]\) is finite.
The factor \([a]\) is finite because \([b]\) is separated.
Explicitly its graph into the fibre product of the finite
composite \([M]\) with the target of \([a]\) is a closed
immersion; the projection of that fibre product is finite
by base change.  Their composite is \([a]\).
This is the source's \srcid{IsFinite.of\_comp} with the
actual endomorphism factor.
Thus \([a]\) is finite.

We retain the source's actual finite-endomorphism flatness
mechanism.  A locally quasi-finite map strictly preserves
strict specialization order: two distinct comparable points
with the same image would be comparable distinct points
of one discrete fibre.  A finite group endomorphism \(f\)
fixes the identity point \(o\).
On the ordered set of generizations of \(o\), it is therefore
strictly monotone.
The generic point \(\eta\) is maximal in that order.
If \(f(\eta)\ne\eta\), its height is strictly less than
that of \(\eta\).  Both heights are finite, being bounded
by the Krull dimension of the Noetherian local ring at \(o\).
But a strictly monotone self-map cannot decrease height,
since it maps each strict chain injectively to a strict chain.
This contradiction proves \(f(\eta)=\eta\).

On an affine neighbourhood of that target generic point,
\(f\) is a finite algebra over a Noetherian domain.
After tensoring with its fraction field choose a finite
vector-space basis.  Clearing denominators in a finite
presentation makes that algebra module free after
localizing at a nonzero element.
Since \(f\) preserves the generic point, the corresponding
source basic open is nonempty.
Thus the source obtains a nonempty open of flat points.
Choose a closed point \(y\) in it; such a point exists
because the finite-type \(k\)-scheme is Jacobson.
As \(k\) is algebraically closed it is a \(k\)-point.
For any closed point \(x\), translation
\(\alpha(z)=z\,y^{-1}x\) on the source and translation by
\(f(y^{-1}x)\) on the target are isomorphisms satisfying
\(f\alpha=\beta f\) and \(\alpha(y)=x\).
Flatness at \(y\), transported through this square of
isomorphisms, proves flatness at \(x\).
Flatness can be tested at closed points in this
Noetherian finite-presentation setting, so \(f\) is flat.
Its image is closed and contains the generic point, hence
is all of the integral \(B\).
This is \srcid{FlatAndSurjectiveOfIsFiniteAbelianEnd}'s
\srcid{apply\_genericPoint},
\srcid{exists\_open\_flat}, and translation argument.
Its rank is locally constant, and \(B\) is connected, so
write \(d(a)\ge1\) for its constant finite-flat rank.
Generic ranks multiply in the function-field tower, giving
\(d(ab)=d(a)d(b)\) for nonzero \(a,b\).

Apply Lemma~\ref{lem:faithful-finite-flat-section-sandwich}
to \([m]\).
The isomorphism \([m]^*H\simeq H^{m^2}\) makes its pulled
section dimensions \(h(m^2t)\).
Lemmas~\ref{lem:faithful-section-growth} and
\ref{lem:faithful-degree-growth-comparison} therefore give
\[
                               d(m)=m^{2g}.
 \tag{DG8}\label{eq:faithful-integer-multiplication-degree}
\]
We have obtained this value by the source's section-growth
argument, rather than taking it as a norm-degree premise.

For an arbitrary finite-flat surjective endomorphism
\(\varphi\) of rank \(\delta\), suppose
\(M h-\varphi^*h\) is effective with \(M\ge1\).
Its nonzero section, raised to the \(t\)-th power, gives an
injection
\[
       H^0(B,\varphi^*H^t)\hookrightarrow H^0(B,H^{Mt}).
\]
Let \(a\) be the shift in
\eqref{eq:faithful-finite-flat-sandwich}.
At \(t=a+1\), that sandwich and
\eqref{eq:faithful-section-growth} give
\[
 \begin{split}
 \delta(a+1)^g
  &\le c\,\delta h(a+1)
   \le c\,h\bigl(M(2a+1)\bigr)\\
  &\le cC M^g(2a+1)^g
   \le cC\,2^g M^g(a+1)^g .
 \end{split}
\]
Cancelling the positive last power yields
\[
                         \delta\le C_0 M^g,\qquad C_0=cC2^g.
 \tag{DG9}\label{eq:faithful-effective-degree-bound}
\]
The shift has disappeared from this constant.

The effective Picard classes form an additive monoid:
products of nonzero sections on the integral \(B\) remain
nonzero.
All \(Q(a)\) are effective because \(H\) is globally
generated.  For each \(a\), some \(Nh-Q(a)\) is effective
by the high-twist section input stated at the start.
Apply Lemma~\ref{lem:faithful-effective-quadratic-bound}
to \(Q\) and an integral basis \(b_i\) of \(O\).
It gives \(C_1\ge1\) with
\[
 C_1T^2h-Q\!\left(\sum_i c_i b_i\right)
                  \quad\text{effective if }|c_i|\le T.
\]
If that element is nonzero, \(T\ge1\).
Equation \eqref{eq:faithful-effective-degree-bound} gives
\[
        d\!\left(\sum_i c_i b_i\right)
                             \le C_0 C_1^g T^{2g}.
\]
Extend \(d\) by any value at \(0\), which none of the
needed inequalities tests.
The hypotheses of
Lemma~\ref{lem:faithful-degree-function-norm} now hold for
\(d\) and \(e=2g\): multiplicativity, positivity,
\eqref{eq:faithful-integer-multiplication-degree}, and this
integral-coordinate bound.
It yields
\(d(a)^{[E:\mathbb Q]}=|N(a)|^{2g}\).

Finally the fibre of \([a]\) over the identity is affine
and finite over \(k\), hence \(\operatorname{Spec}A_a\)
for a finite \(k\)-algebra.
Finite flat base change identifies \(\dim_kA_a=d(a)\),
proving \eqref{eq:faithful-proper-multiplication-norm-degree}.
For an actual germ of the proper datum, its coordinate
dimension equals \(g\): both its tangent points and a
smooth tangent model are bases of the same tangent module.
Thus this is exactly the \(2d'\) exponent used by
\srcid{ExistsFiniteKernelNonzero}.
\end{proof}

\paragraph{Frozen-source ledger.}
The construction-specific flat-cover proof and norm-degree
assembly are in
\href{https://github.com/ChristianKudeba/alpoge-rh-verification/blob/alpoge-source-7cbf7a0c/src/Development/WeilBound/CoreStatements/TorusUniformisation/ChowMahlerFamilies/DegreeMultiplicationNorm.lean}
{\srcid{DegreeMultiplicationNorm}}
(blob \srcid{09a4b78a396ac292ab8fe3938f653e00cca99a6a}):
\srcid{sandwich}, \srcid{finrank\_le},
\srcid{finrank\_act}, \srcid{dim\_eq},
\srcid{exists\_fibre}, and \srcid{finrank\_fibre}.
Its exact wrapper is
\srcid{ExistsFiniteKernelNonzero}
(\srcid{01f7547b58db5fa9282aea4a27d62ae1b3540b55}).
The common denominator, common multiple, and two actual
linear maps are
\srcid{UniformizationProofs.ExistsCommonMultiple};
the chart-to-function-field construction is
\srcid{UniformizationProofs.FinrankSecPullSandwich.chart\_data}.
The latter file's final theorem is phrased for \'etale
covers; the \srcid{sandwich} proof above is the actual
finite-flat extension in \srcid{DegreeMultiplicationNorm}.
It uses neither separability of \(F'/F\) nor an \'etale
assumption.

The filtered-coordinate producers are
\srcid{AbelianSchemeLawProofs.FinrankTotalDegreeImageGrowth},
\srcid{AbelianSchemeLawProofs.FiniteOfNotLinearIndependent}
(\srcid{growth\_of\_veryAmpleBy}), and
\srcid{UniformizationProofs.FiniteMonotoneGrowth}.
The numerical iteration is
\srcid{UniformizationProofs.LeOfPowLePolyMulPow}.
The generic-flat-open and translation proof is
\srcid{AbelianSchemeLawProofs.FlatAndSurjectiveOfIsFiniteAbelianEnd}
(blob \srcid{0a1fb23e45cb9470d7afd3090c9c1e8c3bc80f88}),
not a replacement by a regular/Cohen--Macaulay flatness criterion.
The Serre and ample-twist uses have been identified with
their precise classical statements at the start, so no
source-specific stronger section interface remains implicit.
The actual coordinate-growth, section sandwich, and norm
identification are printed above.
